%% kapa: exam=1 round=30 year=2019 period=1 form=A subject=economics subject_ko=경제학원론 seq=20 number=60
%% source: Q-Net 감정평가사 기출문제 (cst00309) artlSeq=5205205 file=extracted/1008/5205205/2019년 제30회 감정평가사 1차시험 문제지/2019년 제30회 감정평가사 1차시험 문제지/2019년도 제30회 감정평가사 1차 1교시 A형.hwp
%% conversion: kordoc 4.24.51 -> tools/gen1.py (auto); review=proofread-sonnet; flags=none
\begin{question}{60}
하루 24시간을 노동을 하는 시간과 여가를 즐기는 시간으로 양분할 때, 후방굴절형 노동공급곡선이 발생하는 이유는?
\begin{choices}
  \item 임금이 인상될 경우 여가의 가격이 노동의 가격보다 커지기 때문이다.
  \item 임금이 인상될 경우 노동 한 시간 공급으로 할 수 있는 것이 많아지기 때문이다.
  \item 여가가 정상재이고, 소득효과가 대체효과보다 크기 때문이다.
  \item 여가가 정상재이고, 소득효과가 대체효과와 같기 때문이다.
  \item 노동이 열등재이고, 소득효과가 대체효과와 같기 때문이다.
\end{choices}
\end{question}
\answer{3}
